%!TEX root = ../main.tex
%----------------------------------------------------------------------
%  Section: The Theory-Informed KRR Estimator
%  Paper:   Theory-Informed Kernel Ridge Regression for Asset Pricing
%  Authors: Amedeo Andriollo and Juan F. Imbet
%----------------------------------------------------------------------

\section{The Theory-Informed Kernel Ridge Regression Estimator}
\label{sec:framework}

This section develops the econometric framework.
Section~\ref{sub:setup} introduces notation and the prediction environment.
Section~\ref{sub:estimator} defines the theory-informed KRR objective.
Sections~\ref{sub:penalties}--\ref{sub:finite_dim} formalize the structural penalty types
and establish the representer theorem and finite-dimensional reduction that
make the estimator computationally tractable.
Section~\ref{sub:nesting} shows that the framework nests several influential
estimators as special cases.
Sections~\ref{sub:cv}--\ref{sub:inference} address implementation, misspecification
robustness, and inference.

%======================================================================
\subsection{Setup and Notation}
\label{sub:setup}
%======================================================================

We observe a panel of $N$ assets over $T$ periods.
Let $R_{i,t+1}$ denote the excess return of asset $i$ realized at date $t+1$,
and let $\mathbf{X}_{i,t} \in \mathbb{R}^p$ collect firm-level characteristics
and macroeconomic state variables observable at date $t$.
The total sample size is $n = NT$.

\begin{assumption}[Data-generating process]
\label{ass:dgp}
There exists a measurable function $f^{*}: \mathbb{R}^{p} \to \mathbb{R}$
such that
\begin{equation}
\label{eq:dgp}
R_{i,t+1} = f^{*}(\mathbf{X}_{i,t}) + \varepsilon_{i,t+1},
\qquad
\mathbb{E}[\varepsilon_{i,t+1} \mid \mathbf{X}_{i,t}] = 0,
\end{equation}
where $f^{*}(\mathbf{X}_{i,t}) = \mathbb{E}[R_{i,t+1} \mid \mathbf{X}_{i,t}]$
is the conditional expected excess return.
\end{assumption}

We approximate $f^{*}$ within a reproducing kernel Hilbert space (RKHS).

\begin{definition}[RKHS]
\label{def:rkhs}
Let $k: \mathbb{R}^p \times \mathbb{R}^p \to \mathbb{R}$ be a symmetric,
positive-definite kernel.
The RKHS $\mathcal{H}$ associated with $k$ is the completion of
$\operatorname{span}\{k(\mathbf{x}, \cdot) : \mathbf{x} \in \mathbb{R}^p\}$
under the inner product satisfying the reproducing property
$f(\mathbf{x}) = \langle f, k(\mathbf{x}, \cdot) \rangle_{\mathcal{H}}$
for all $f \in \mathcal{H}$.
The RKHS norm $\|f\|_{\mathcal{H}}^2 = \langle f, f \rangle_{\mathcal{H}}$
measures the complexity of $f$.
\end{definition}

We stack observations using the index $s = (i,t)$, so that
$\mathbf{X}_s \equiv \mathbf{X}_{i,t}$ and $R_s \equiv R_{i,t+1}$ for
$s = 1, \ldots, n$.
The $n \times n$ kernel (Gram) matrix is
$\mathbf{K}$, with $[\mathbf{K}]_{s,s'} = k(\mathbf{X}_s, \mathbf{X}_{s'})$.
The return vector is $\mathbf{R} = (R_1, \ldots, R_n)^\top \in \mathbb{R}^n$.

%======================================================================
\subsection{The Estimator}
\label{sub:estimator}
%======================================================================

Our estimator augments the standard KRR objective with a collection of
structural penalty terms, each encoding restrictions implied by a specific
economic model.

\begin{definition}[Theory-Informed KRR]
\label{def:tikrr}
Let $\{C_j\}_{j=1}^{J}$ be a collection of structural penalty functionals
$C_j: \mathcal{H} \to \mathbb{R}_+$,
where each $C_j(f) \geq 0$ with equality when $f$ satisfies the
restrictions of economic model~$j$.
Let $\lambda_{\mathrm{stat}} > 0$ and $\mu_j \geq 0$, $j = 1, \ldots, J$.
The \emph{theory-informed KRR} estimator is
\begin{equation}
\label{eq:tikrr}
\hat{f} \;=\; \arg\min_{f \in \mathcal{H}} \;
\underbrace{\frac{1}{n}\sum_{s=1}^{n}
  \bigl(R_s - f(\mathbf{X}_s)\bigr)^2}_{\text{prediction loss}}
\;+\;
\underbrace{\lambda_{\mathrm{stat}}\,\|f\|_{\mathcal{H}}^2}_{\text{statistical regularization}}
\;+\;
\underbrace{\sum_{j=1}^{J} \mu_j\, C_j(f)}_{\text{structural penalties}}.
\end{equation}
\end{definition}

The first two terms constitute standard kernel ridge regression:
the squared-error prediction loss trades off against the RKHS-norm penalty,
which controls the smoothness of $f$.
The third term is new.
Each $C_j$ penalizes deviations from a specific economic model's
predictions, and the associated multiplier $\mu_j$ governs how strongly
that model's restrictions influence the estimate.
When $\mu_j = 0$, model $j$ is ignored; as $\mu_j \to \infty$,
the restrictions of model $j$ are imposed exactly.

\begin{remark}[Interpretation of multipliers]
\label{rem:multipliers}
The vector $(\lambda_{\mathrm{stat}}, \mu_1, \ldots, \mu_J)$ balances
agnostic statistical regularization against theory-informed regularization.
A large $\lambda_{\mathrm{stat}}$ relative to $\mu_j$ yields an estimate
close to standard KRR; large $\mu_j$ relative to $\lambda_{\mathrm{stat}}$
forces the estimator to respect model $j$'s restrictions even if doing
so worsens in-sample fit.
Cross-validation selects the balance that maximizes out-of-sample
predictive accuracy (Section~\ref{sub:cv}).
\end{remark}

%======================================================================
\subsection{Structural Penalty Types}
\label{sub:penalties}
%======================================================================

We formalize four classes of structural penalties.
Each encodes restrictions from a distinct family of asset-pricing models.
A key requirement---verified explicitly below---is that every $C_j(f)$
depends on $f$ only through its evaluations at the observed data points
$\{f(\mathbf{X}_s)\}_{s=1}^n$.

%----------------------------------------------------------------------
\subsubsection{Type~A: Euler Equation Penalties}
\label{sub:euler}
%----------------------------------------------------------------------

Let model $j$ specify a stochastic discount factor (SDF)
$M_{j,t+1} = M_j(\theta_j, \mathbf{X}_t, \mathbf{X}_{t+1})$
parameterized by $\theta_j$, with a preliminary estimator
$\hat{\theta}_j$ obtained by GMM or maximum likelihood.
The Euler equation requires
$\mathbb{E}_t[M_{j,t+1}\, R_{i,t+1}^{\mathrm{gross}}] = 1$
for every asset $i$ and date $t$,
where $R_{i,t+1}^{\mathrm{gross}} = f(\mathbf{X}_{i,t}) + R_t^f$
is the gross return and $R_t^f$ is the risk-free rate.

\begin{definition}[Euler equation penalty]
\label{def:euler_penalty}
For SDF model $j$ with pre-estimated parameter $\hat{\theta}_j$, define
\begin{equation}
\label{eq:euler_penalty}
C_j^{\mathrm{Euler}}(f) \;=\;
\sum_{t=1}^{T}
\left\|
\frac{1}{N}\sum_{i=1}^{N}
M_j\bigl(\hat{\theta}_j, \mathbf{X}_t, \mathbf{X}_{t+1}\bigr)
\cdot
\bigl(f(\mathbf{X}_{i,t}) + R_t^f\bigr) - 1
\right\|^2.
\end{equation}
\end{definition}

\noindent
The penalty is zero precisely when the cross-sectional Euler equation
holds at every date $t$ given $\hat{\theta}_j$.
Since $M_j(\hat{\theta}_j, \mathbf{X}_t, \mathbf{X}_{t+1})$ and $R_t^f$ are data,
the penalty depends on $f$ only through $\{f(\mathbf{X}_{i,t})\}_{i,t}$.

\begin{remark}[Coverage]
Type~A penalties accommodate consumption-based models (CCAPM, long-run risk,
rare disasters, habit formation), as well as production-based SDFs derived
from firms' intertemporal optimality conditions.
Each distinct SDF specification generates a separate penalty $C_j^{\mathrm{Euler}}$.
\end{remark}

%----------------------------------------------------------------------
\subsubsection{Type~B: Production and Investment FOC Penalties}
\label{sub:production}
%----------------------------------------------------------------------

Production-based asset pricing relates expected returns to observable
investment and profitability variables through firms' first-order conditions.
Let model $j$ imply a mapping $g_j: \mathbb{R}^{d_j} \to \mathbb{R}$
from investment characteristics $\mathbf{Z}_{i,t}^{(j)}$
(e.g., investment-to-capital ratio $I_{i,t}/K_{i,t}$, return on equity
$\mathrm{ROE}_{i,t}$, asset growth) to expected returns.

\begin{definition}[Production FOC penalty]
\label{def:prod_penalty}
For production model $j$ with predicted returns $g_j(\mathbf{Z}_{i,t}^{(j)})$, define
\begin{equation}
\label{eq:prod_penalty}
C_j^{\mathrm{prod}}(f) \;=\;
\frac{1}{n}\sum_{s=1}^{n}
\bigl(
f(\mathbf{X}_s) - g_j\bigl(\mathbf{Z}_s^{(j)}\bigr)
\bigr)^2.
\end{equation}
\end{definition}

\noindent
The penalty is a mean-squared deviation between the nonparametric estimate
and the production model's implied expected return.
Since $g_j(\mathbf{Z}_s^{(j)})$ is computed from data, the penalty depends on $f$
only through $\{f(\mathbf{X}_s)\}_{s=1}^n$.

\begin{remark}[Nesting the $q$-theory framework]
When $g_j$ encodes the first-order condition from the investment CAPM
or the $q$-factor model,
Type~B penalties embed the production approach
to asset pricing directly into the nonparametric objective.
\end{remark}

%----------------------------------------------------------------------
\subsubsection{Type~C: Factor Structure Penalties}
\label{sub:factor}
%----------------------------------------------------------------------

Factor models impose that expected returns obey a linear pricing
relation with respect to exposure to priced factors.
Let model $j$ specify a factor $F_{j,t}$ (e.g., an intermediary capital
ratio factor, market excess return, or a latent factor).
Let $\beta_{i,j}$ denote asset $i$'s loading on $F_{j,t}$,
pre-estimated by time-series regression, and let $\alpha_{i,j}$
be the corresponding intercept.

\begin{definition}[Factor structure penalty]
\label{def:factor_penalty}
For factor model $j$ with estimated loadings
$\hat{\beta}_{i,j}$, $\hat{\alpha}_{i,j}$ ($i=1,\ldots,N$)
and factor realizations $F_{j,t}$ ($t=1,\ldots,T$), define
\begin{equation}
\label{eq:factor_penalty}
C_j^{\mathrm{factor}}(f) \;=\;
\frac{1}{n}\sum_{i=1}^{N}\sum_{t=1}^{T}
\bigl(
f(\mathbf{X}_{i,t}) - \hat{\beta}_{i,j}\, F_{j,t} - \hat{\alpha}_{i,j}
\bigr)^2.
\end{equation}
\end{definition}

\noindent
When asset pricing restricts $\hat{\alpha}_{i,j} = 0$ for all $i$
(exact factor pricing), the penalty reduces to
$\frac{1}{n}\sum_{i,t}(f(\mathbf{X}_{i,t}) - \hat{\beta}_{i,j}\, F_{j,t})^2$,
penalizing deviations from the factor model's predictions of expected returns.
Again, the penalty depends on $f$ only through $\{f(\mathbf{X}_{i,t})\}_{i,t}$.

\begin{remark}[Intermediary asset pricing]
When $F_{j,t}$ is the intermediary capital ratio or the intermediary leverage factor,
Type~C penalties embed intermediary-based
models into the nonparametric estimator.
This is particularly valuable when the intermediary model captures
time-variation in risk premia that firm-level characteristics alone cannot.
\end{remark}

%----------------------------------------------------------------------
\subsubsection{Type~D: Demand-System Equilibrium Penalties}
\label{sub:demand}
%----------------------------------------------------------------------

Demand-system approaches to asset pricing derive equilibrium prices
from investors' portfolio demand functions.
Let $D_i: \mathbb{R}^N \to \mathbb{R}^N$ map the vector of expected returns to investor~$i$'s
optimal portfolio weights, and let $w_i > 0$ denote
investor~$i$'s wealth share, with the index now running over investors
rather than assets.
Let $\mathbf{S}_t \in \mathbb{R}^N$ denote the vector of asset market-capitalization
shares (the supply side) at date $t$.

\begin{definition}[Demand-system equilibrium penalty]
\label{def:demand_penalty}
For demand-system model $j$ with demand functions
$\{D_i^{(j)}\}$ and observed supplies $\{\mathbf{S}_t\}$, let
$\hat{\mathbf{f}}_t = (f(\mathbf{X}_{1,t}), \ldots, f(\mathbf{X}_{N,t}))^\top$
collect expected-return evaluations.  Define
\begin{equation}
\label{eq:demand_penalty}
C_j^{\mathrm{demand}}(f) \;=\;
\sum_{t=1}^{T}
\left\|
\sum_{i} w_i^{(j)}\, D_i^{(j)}\bigl(\hat{\mathbf{f}}_t\bigr) - \mathbf{S}_t
\right\|^2.
\end{equation}
\end{definition}

\noindent
The penalty measures the squared market-clearing failure:
it is zero when the demand system clears at the predicted expected returns.
It depends on $f$ only through $\{f(\mathbf{X}_{i,t})\}_{i,t}$.

\begin{remark}[Connections to the demand-system literature]
The demand-system approach models institutional
investors' demand as a function of characteristics and expected returns.
Type~D penalties encode the equilibrium restriction that aggregate demand
equals supply, translating the demand-system estimator's logic into a
penalty within our nonparametric framework.
When $D_i^{(j)}$ is linear in expected returns, $C_j^{\mathrm{demand}}$
is quadratic in $\hat{\mathbf{f}}_t$ and the overall problem remains convex.
\end{remark}

%======================================================================
\subsection{The Representer Theorem}
\label{sub:representer}
%======================================================================

The following result establishes that the infinite-dimensional
optimization problem~\eqref{eq:tikrr} admits a finite-dimensional
solution, despite the presence of the structural penalties.

\begin{proposition}[Representer theorem for theory-informed KRR]
\label{prop:representer}
Let $\mathcal{H}$ be an RKHS with kernel $k$.
Suppose each structural penalty $C_j: \mathcal{H} \to \mathbb{R}_+$
depends on $f$ only through the vector of evaluations
$(f(\mathbf{X}_1), \ldots, f(\mathbf{X}_n))^\top \in \mathbb{R}^n$.
That is, there exist functions $\tilde{C}_j: \mathbb{R}^n \to \mathbb{R}_+$
such that
\begin{equation}
\label{eq:penalty_eval}
C_j(f) = \tilde{C}_j\bigl(f(\mathbf{X}_1), \ldots, f(\mathbf{X}_n)\bigr)
\qquad \text{for all } f \in \mathcal{H}, \; j = 1, \ldots, J.
\end{equation}
Then any minimizer of~\eqref{eq:tikrr} admits the representation
\begin{equation}
\label{eq:representer}
\hat{f}(\cdot) = \sum_{s=1}^{n} \hat{\alpha}_s\, k(\mathbf{X}_s, \cdot),
\end{equation}
for some $\hat{\boldsymbol{\alpha}} = (\hat{\alpha}_1, \ldots, \hat{\alpha}_n)^\top
\in \mathbb{R}^n$.
\end{proposition}

\begin{proof}
\emph{The proof is given in Appendix~\ref{app:proofs}.}
\end{proof}

\begin{remark}[Verification for Types~A--D]
\label{rem:verify_representer}
The condition~\eqref{eq:penalty_eval} holds for all four penalty types:
\begin{enumerate}[label=(\alph*)]
\item \emph{Type~A (Euler):}
  $C_j^{\mathrm{Euler}}(f)$ depends on $f$ through
  $\{f(\mathbf{X}_{i,t})\}_{i,t}$, since
  $M_j(\hat{\theta}_j, \mathbf{X}_t, \mathbf{X}_{t+1})$ and $R_t^f$
  are data.
\item \emph{Type~B (Production):}
  $C_j^{\mathrm{prod}}(f)$ is explicitly a function of
  $\{f(\mathbf{X}_s)\}_{s}$ and data $\{g_j(\mathbf{Z}_s^{(j)})\}_s$.
\item \emph{Type~C (Factor):}
  $C_j^{\mathrm{factor}}(f)$ involves $\{f(\mathbf{X}_{i,t})\}_{i,t}$
  and pre-estimated $\hat{\beta}_{i,j}$, $\hat{\alpha}_{i,j}$, $F_{j,t}$.
\item \emph{Type~D (Demand):}
  $C_j^{\mathrm{demand}}(f)$ depends on $f$ through
  $\hat{\mathbf{f}}_t = (f(\mathbf{X}_{1,t}), \ldots, f(\mathbf{X}_{N,t}))^\top$.
\end{enumerate}
In each case, the penalty is a function of the $n$-dimensional evaluation
vector, so Proposition~\ref{prop:representer} applies.
This extends classical representer theorems to accommodate theory-informed
structural penalties.
\end{remark}

%======================================================================
\subsection{Finite-Dimensional Reduction}
\label{sub:finite_dim}
%======================================================================

By Proposition~\ref{prop:representer}, it suffices to optimize
over $\boldsymbol{\alpha} \in \mathbb{R}^n$.
Substituting $f(\cdot) = \sum_s \alpha_s\, k(\mathbf{X}_s, \cdot)$ into
the objective yields the following.

\begin{proposition}[Finite-dimensional problem]
\label{prop:finite_dim}
Under the representer theorem (Proposition~\ref{prop:representer}),
the theory-informed KRR estimator~\eqref{eq:tikrr} reduces to
\begin{equation}
\label{eq:finite_dim}
\hat{\boldsymbol{\alpha}} \;=\; \arg\min_{\boldsymbol{\alpha} \in \mathbb{R}^n} \;
\frac{1}{n}\|\mathbf{R} - \mathbf{K}\boldsymbol{\alpha}\|^2
\;+\;
\lambda_{\mathrm{stat}}\,\boldsymbol{\alpha}^\top \mathbf{K}\, \boldsymbol{\alpha}
\;+\;
\sum_{j=1}^{J} \mu_j\, \tilde{C}_j(\mathbf{K}\boldsymbol{\alpha}),
\end{equation}
where $\mathbf{K}$ is the $n \times n$ kernel matrix,
$\hat{\mathbf{f}} = \mathbf{K}\boldsymbol{\alpha}$ is the $n \times 1$
vector of fitted values, and
$\tilde{C}_j(\hat{\mathbf{f}}) = C_j(f)$ expresses each penalty
in terms of fitted values.
\end{proposition}

\begin{proof}
\emph{The proof is given in Appendix~\ref{app:proofs}.}
\end{proof}

%----------------------------------------------------------------------
\subsubsection{Convexity and First-Order Conditions for Quadratic Penalties}
\label{sub:foc}
%----------------------------------------------------------------------

When all active structural penalties are quadratic in the fitted values
$\hat{\mathbf{f}} = \mathbf{K}\boldsymbol{\alpha}$, the finite-dimensional
problem~\eqref{eq:finite_dim} becomes a convex quadratic program with
a closed-form solution.

For Type~A penalties, define the $T \times n$ matrix $\mathbf{A}_j$
with rows indexed by $t$:
\begin{equation}
\label{eq:A_matrix}
[\mathbf{A}_j]_{t,s} =
\begin{cases}
\dfrac{1}{N}\, M_j\bigl(\hat{\theta}_j, \mathbf{X}_t, \mathbf{X}_{t+1}\bigr)
& \text{if } s = (i,t) \text{ for some } i, \\[4pt]
0 & \text{otherwise},
\end{cases}
\end{equation}
and the $T \times 1$ vector $\mathbf{b}_j$ with
$[\mathbf{b}_j]_t = 1 - \frac{1}{N}\sum_i M_j(\hat{\theta}_j,
\mathbf{X}_t, \mathbf{X}_{t+1})\, R_t^f$.
Then
\begin{equation}
\label{eq:euler_quad}
C_j^{\mathrm{Euler}}(f) =
\|\mathbf{A}_j\, \hat{\mathbf{f}} - \mathbf{b}_j\|^2
= (\mathbf{A}_j\, \mathbf{K}\boldsymbol{\alpha} - \mathbf{b}_j)^\top
  (\mathbf{A}_j\, \mathbf{K}\boldsymbol{\alpha} - \mathbf{b}_j).
\end{equation}

For Type~B penalties, let $\mathbf{g}_j \in \mathbb{R}^n$ with
$[\mathbf{g}_j]_s = g_j(\mathbf{Z}_s^{(j)})$.  Then
\begin{equation}
\label{eq:prod_quad}
C_j^{\mathrm{prod}}(f) =
\frac{1}{n}\|\mathbf{K}\boldsymbol{\alpha} - \mathbf{g}_j\|^2.
\end{equation}

For Type~C penalties, let $\mathbf{h}_j \in \mathbb{R}^n$ with
$[\mathbf{h}_j]_s = \hat{\beta}_{i,j}\, F_{j,t} + \hat{\alpha}_{i,j}$
where $s = (i,t)$.  Then
\begin{equation}
\label{eq:factor_quad}
C_j^{\mathrm{factor}}(f) =
\frac{1}{n}\|\mathbf{K}\boldsymbol{\alpha} - \mathbf{h}_j\|^2.
\end{equation}

\begin{proposition}[First-order conditions]
\label{prop:foc}
Suppose every active penalty ($\mu_j > 0$) is of Type~A, B, or C
(or more generally, quadratic in $\hat{\mathbf{f}}$).
Define the aggregated penalty matrix
\begin{equation}
\label{eq:Omega}
\boldsymbol{\Omega} \;=\;
\sum_{j \in \mathcal{J}^A} \mu_j\, \mathbf{A}_j^\top \mathbf{A}_j
\;+\;
\sum_{j \in \mathcal{J}^B} \frac{\mu_j}{n}\, \mathbf{I}_n
\;+\;
\sum_{j \in \mathcal{J}^C} \frac{\mu_j}{n}\, \mathbf{I}_n,
\end{equation}
where $\mathcal{J}^A$, $\mathcal{J}^B$, $\mathcal{J}^C$ index
Type~A, B, and C penalties respectively, and the aggregated target vector
\begin{equation}
\label{eq:omega_vec}
\boldsymbol{\omega} \;=\;
\sum_{j \in \mathcal{J}^A} \mu_j\, \mathbf{A}_j^\top \mathbf{b}_j
\;+\;
\sum_{j \in \mathcal{J}^B} \frac{\mu_j}{n}\, \mathbf{g}_j
\;+\;
\sum_{j \in \mathcal{J}^C} \frac{\mu_j}{n}\, \mathbf{h}_j.
\end{equation}
Then the first-order conditions for~\eqref{eq:finite_dim} are
\begin{equation}
\label{eq:foc}
\left[
\frac{1}{n}\,\mathbf{K}
+ \lambda_{\mathrm{stat}}\,\mathbf{I}_n
+ \boldsymbol{\Omega}\,\mathbf{K}
\right] \boldsymbol{\alpha}
\;=\;
\frac{1}{n}\,\mathbf{R}
+ \boldsymbol{\omega},
\end{equation}
and when $\mathbf{K}$ is strictly positive definite, the unique solution is
\begin{equation}
\label{eq:alpha_star}
\hat{\boldsymbol{\alpha}} \;=\;
\left[
\frac{1}{n}\,\mathbf{K}
+ \lambda_{\mathrm{stat}}\,\mathbf{I}_n
+ \boldsymbol{\Omega}\,\mathbf{K}
\right]^{-1}
\left(
\frac{1}{n}\,\mathbf{R} + \boldsymbol{\omega}
\right).
\end{equation}
\end{proposition}

\begin{proof}
\emph{The proof is given in Appendix~\ref{app:proofs}.}
\end{proof}

\begin{remark}[Computational cost]
\label{rem:computation}
The dominant cost is solving the $n \times n$ linear system~\eqref{eq:alpha_star},
which requires $O(n^3)$ operations.
When $n$ is large, one can exploit the eigendecomposition
$\mathbf{K} = \mathbf{U}\boldsymbol{\Lambda}\mathbf{U}^\top$
(computed once at cost $O(n^3)$) and solve in the rotated basis,
where the matrix becomes diagonal plus a low-rank perturbation from
the Type~A penalties.
Alternatively, random Fourier features or Nystr\"{o}m approximations
reduce the effective dimension.
\end{remark}

\begin{remark}[Type~D penalties and convexity]
\label{rem:type_d_convexity}
When demand functions $D_i^{(j)}$ are linear in expected returns,
$C_j^{\mathrm{demand}}$ is quadratic in $\hat{\mathbf{f}}$ and the
problem remains a convex QP with a structure analogous to~\eqref{eq:foc}.
For nonlinear demand systems, the problem is generally nonconvex and
requires iterative methods (e.g., projected gradient descent or
alternating minimization), though convexity of the remaining terms
provides a useful warm start.
\end{remark}

%======================================================================
\subsection{Nesting Structure}
\label{sub:nesting}
%======================================================================

A key advantage of the framework is that it nests several influential
estimators as special or limiting cases.
This both clarifies the relationship between existing methods and
suggests that the theory-informed KRR estimator can weakly dominate
each special case.

\begin{proposition}[Nesting]
\label{prop:nesting}
The theory-informed KRR estimator~\eqref{eq:tikrr} nests the following:
\begin{enumerate}[label=(\roman*)]
\item \textbf{Standard KRR.}
  Setting $\mu_j = 0$ for all $j = 1, \ldots, J$ yields
  \[
  \hat{f}^{\mathrm{KRR}} = \arg\min_{f \in \mathcal{H}}
  \frac{1}{n}\sum_s (R_s - f(\mathbf{X}_s))^2
  + \lambda_{\mathrm{stat}}\|f\|_{\mathcal{H}}^2,
  \]
  the standard kernel ridge regression estimator.

\item \textbf{Constrained minimum-distance / GMM.}
  Setting $\lambda_{\mathrm{stat}} = 0$ and activating a single
  Type~A penalty with $\mu_j \to \infty$ yields
  \[
  \hat{f}^{\mathrm{GMM}} = \arg\min_{f \in \mathcal{H}}
  \frac{1}{n}\sum_s (R_s - f(\mathbf{X}_s))^2
  \quad \text{s.t.} \quad
  C_j^{\mathrm{Euler}}(f) = 0,
  \]
  which is a nonparametric generalization of GMM-based SDF estimation.
  When $\mathcal{H}$ is restricted to parametric functions,
  this recovers classical GMM.

\item \textbf{Approximate instrumented PCA.}
  Consider a linear kernel $k(\mathbf{x}, \mathbf{x}') = \mathbf{x}^\top \mathbf{x}'$
  so that $f(\mathbf{X}_{i,t}) = \boldsymbol{\beta}^\top \mathbf{X}_{i,t}$
  for some $\boldsymbol{\beta} \in \mathbb{R}^p$, and a Type~C penalty
  where the factor $F_{j,t}$ is a target-PCA factor.
  Then~\eqref{eq:tikrr} becomes
  \begin{equation}
  \label{eq:ipca_nest}
  \hat{\boldsymbol{\beta}} = \arg\min_{\boldsymbol{\beta}}
  \frac{1}{n}\sum_{i,t}(R_{i,t+1} - \boldsymbol{\beta}^\top \mathbf{X}_{i,t})^2
  + \lambda_{\mathrm{stat}}\|\boldsymbol{\beta}\|^2
  + \mu_j \frac{1}{n}\sum_{i,t}
    (\boldsymbol{\beta}^\top \mathbf{X}_{i,t}
     - \hat{\beta}_{i,j} F_{j,t})^2,
  \end{equation}
  a ridge regression with an additional penalty encouraging alignment
  with the target-PCA factor structure.
  This approximates the cross-sectional
  and time-series restrictions in instrumented PCA methods.

\item \textbf{Text-informed penalized regression.}
  Replacing $\mathcal{H}$ with a finite-dimensional linear space
  $\mathbb{R}^p$, substituting the RKHS-norm penalty with an
  $\ell_1$ norm $\|\boldsymbol{\beta}\|_1$, and setting
  $C_j(\boldsymbol{\beta}) = \sum_{k=1}^p \gamma_{j,k}\, |\beta_k|$
  where $\gamma_{j,k}$ are text-derived importance weights, yields
  \begin{equation}
  \label{eq:llm_lasso_nest}
  \hat{\boldsymbol{\beta}} = \arg\min_{\boldsymbol{\beta}}
  \frac{1}{n}\sum_s (R_s - \boldsymbol{\beta}^\top \mathbf{X}_s)^2
  + \lambda_{\mathrm{stat}} \|\boldsymbol{\beta}\|_1
  + \mu_j \sum_k \gamma_{j,k}\, |\beta_k|,
  \end{equation}
  a variant of the adaptive Lasso with text-informed variable-specific
  penalties, nesting the LLM-Lasso approach of
  incorporating domain knowledge from large language models into
  penalized regression.
\end{enumerate}
\end{proposition}

\begin{proof}
\emph{The proof is given in Appendix~\ref{app:proofs}.}
\end{proof}

\begin{remark}[Weak dominance]
\label{rem:weak_dominance}
Since standard KRR corresponds to $\mu_j = 0$ for all $j$,
and the cross-validated theory-informed KRR is free to select
$\hat{\mu}_j = 0$ if theory $j$ is unhelpful, the theory-informed
estimator weakly dominates standard KRR in out-of-sample
prediction loss---up to the granularity of the hyperparameter search.
The same logic applies to each nested special case: the larger
hyperparameter space can only improve (or match) the best
achievable out-of-sample performance.
\end{remark}

%======================================================================
\subsection{Cross-Validation for Hyperparameter Selection}
\label{sub:cv}
%======================================================================

The hyperparameter vector
$\boldsymbol{\lambda} = (\lambda_{\mathrm{stat}}, \mu_1, \ldots, \mu_J)$
is $(J+1)$-dimensional.
We introduce grouped penalties and temporal cross-validation
to make the selection tractable.

\paragraph{Grouped penalties.}
We assign each model $j$ to one of $G$ groups based on its
economic type.  Let $g: \{1, \ldots, J\} \to \{1, \ldots, G\}$
denote the group assignment, and impose $\mu_j = \mu_{g(j)}$
for all $j$.  In our empirical application, $G = 8$, corresponding to
the four penalty types (A--D) each potentially with two variants
(e.g., consumption-based vs.\ production-based SDFs within Type~A).
This reduces the hyperparameter space from $(J+1)$ to $(G+1) = 9$ dimensions.

\paragraph{Temporal cross-validation.}
Financial panel data exhibit strong serial dependence, so we use
leave-one-year-out (LOYO) cross-validation to select $\boldsymbol{\lambda}$.
Partition the sample into $T_{\mathrm{yr}}$ calendar years
$\mathcal{T}_1, \ldots, \mathcal{T}_{T_{\mathrm{yr}}}$.
For fold $\ell$, let
$\mathcal{S}_{-\ell} = \{(i,t) : t \notin \mathcal{T}_\ell\}$ be the
training set and $\mathcal{S}_\ell = \{(i,t) : t \in \mathcal{T}_\ell\}$
the validation set.
The cross-validated objective is
\begin{equation}
\label{eq:cv}
\mathrm{CV}(\boldsymbol{\lambda}) \;=\;
\frac{1}{T_{\mathrm{yr}}} \sum_{\ell=1}^{T_{\mathrm{yr}}}
\frac{1}{|\mathcal{S}_\ell|}
\sum_{s \in \mathcal{S}_\ell}
\bigl(R_s - \hat{f}_{-\ell}(\mathbf{X}_s; \boldsymbol{\lambda})\bigr)^2,
\end{equation}
where $\hat{f}_{-\ell}(\cdot; \boldsymbol{\lambda})$ is the theory-informed
KRR estimator trained on $\mathcal{S}_{-\ell}$ with hyperparameters
$\boldsymbol{\lambda}$.

\paragraph{Search strategy.}
We minimize $\mathrm{CV}(\boldsymbol{\lambda})$ by random search over
a 9-dimensional log-uniform grid:
$\log_{10}\lambda_{\mathrm{stat}} \in [-6, 2]$ and
$\log_{10}\mu_g \in [-6, 2]$ for $g = 1, \ldots, G$.
We draw $B = 500$ random configurations and evaluate each by solving
the finite-dimensional problem~\eqref{eq:finite_dim} $T_{\mathrm{yr}}$
times (once per fold).
The selected hyperparameters are
$\hat{\boldsymbol{\lambda}} = \arg\min_{\boldsymbol{\lambda}}
\mathrm{CV}(\boldsymbol{\lambda})$.

\begin{remark}[Computational considerations]
\label{rem:cv_computation}
Each fold requires solving a linear system of size $|\mathcal{S}_{-\ell}|$,
which is $O(((T_{\mathrm{yr}}-1) \cdot N / T_{\mathrm{yr}})^3)$ per fold.
With $T_{\mathrm{yr}} \approx 30$--$60$ years and $N \approx 3{,}000$--$10{,}000$
assets, this is demanding but feasible with precomputed kernel
eigendecompositions and parallelization across folds and random draws.
The total cost scales as $O(B \cdot T_{\mathrm{yr}} \cdot n_{\mathrm{fold}}^3)$
where $n_{\mathrm{fold}} = |\mathcal{S}_{-\ell}|$.
\end{remark}

%======================================================================
\subsection{Misspecification Robustness}
\label{sub:misspecification}
%======================================================================

A central concern with theory-informed estimation is that imposing
restrictions from a misspecified model may worsen predictions.
The following result shows that cross-validation provides an automatic
safeguard.

\begin{proposition}[Misspecification robustness]
\label{prop:misspecification}
Suppose model $j$ is misspecified in the sense that the true
conditional expected return function $f^*$ satisfies
$C_j(f^*) > 0$.
Under the following regularity conditions:
\begin{enumerate}[label=\textnormal{(R\arabic*)}]
\item \label{cond:consistency}
  The RKHS $\mathcal{H}$ is rich enough that
  $\inf_{f \in \mathcal{H}} \|f - f^*\|_{L^2} = 0$;
\item \label{cond:cv_oracle}
  The cross-validation procedure~\eqref{eq:cv} is asymptotically
  oracle-efficient in the sense that
  $\mathrm{CV}(\hat{\boldsymbol{\lambda}}) - \inf_{\boldsymbol{\lambda}}
  \mathrm{CV}(\boldsymbol{\lambda}) = o_p(1)$;
\item \label{cond:bounded}
  The penalty $C_j$ is continuous in $\|\cdot\|_{L^2}$ and the
  parameter space for $\mu_j$ includes zero;
\end{enumerate}
the cross-validated multiplier satisfies
$\hat{\mu}_j \xrightarrow{p} 0$ as $n \to \infty$.
\end{proposition}

\begin{proof}
\emph{The proof is given in Appendix~\ref{app:proofs}.}
\end{proof}

\begin{remark}[Practical implication]
\label{rem:misspec_practical}
Proposition~\ref{prop:misspecification} provides a theoretical guarantee
that misspecified theories ``turn themselves off'' through cross-validation.
In finite samples, a moderately misspecified theory may still receive
positive $\hat{\mu}_j$ if it captures useful structure despite not
being exactly correct.
The estimator thus exploits partial truths---a theory need not be
exactly right to be useful for prediction.
\end{remark}

\begin{remark}[Comparison to Bayesian model averaging]
\label{rem:bma}
The misspecification robustness of theory-informed KRR is analogous to
Bayesian model averaging (BMA) in spirit: both frameworks downweight
models that fit poorly out of sample.
However, the mechanism differs.
BMA assigns posterior probabilities to discrete model specifications,
whereas theory-informed KRR treats each model's restrictions as a
continuous penalty with cross-validated strength.
The latter avoids the combinatorial explosion of the BMA model space
and allows partial incorporation of each theory's restrictions.
\end{remark}

%======================================================================
\subsection{Inference on Model Value}
\label{sub:inference}
%======================================================================

Beyond prediction, one may wish to test whether a particular economic
model contributes statistically significant information beyond agnostic
regularization.
This amounts to testing
\begin{equation}
\label{eq:test_mu}
H_0: \mu_j^* = 0
\qquad \text{vs.} \qquad
H_1: \mu_j^* > 0,
\end{equation}
where $\mu_j^*$ is the population-optimal multiplier for model $j$.
Rejection of $H_0$ implies that model $j$ provides marginal predictive
value beyond what the kernel and other models already capture.

%----------------------------------------------------------------------
\subsubsection{Connection to Uniform Inference for KRR}
\label{sub:uniform_inference}
%----------------------------------------------------------------------

Recent work by \citet*{SinghVijaykumar2023} establishes uniform
confidence bands for kernel ridge regression estimators of conditional
expectations.
Their framework provides pointwise and uniform inference for
$f^*(\mathbf{x})$ across $\mathbf{x} \in \mathcal{X}$ using
a bias-corrected estimator and variance estimation based on the
kernel's spectral properties.

We outline how this can be extended to test~\eqref{eq:test_mu}.

\begin{proposition}[Testing the marginal value of economic models]
\label{prop:test_mu}
Suppose the conditions of \citet*{SinghVijaykumar2023} hold for
the RKHS $\mathcal{H}$, augmented with the following:
\begin{enumerate}[label=\textnormal{(T\arabic*)}]
\item \label{cond:test_smooth}
  The mapping $\mu_j \mapsto \hat{f}(\cdot\,; \mu_j)$ is differentiable
  in $L^2$, with derivative $\partial_{\mu_j} \hat{f}$ satisfying
  $\|\partial_{\mu_j} \hat{f}\|_{L^2} = O(1)$;
\item \label{cond:test_rate}
  The estimation error satisfies
  $\|\hat{f}(\cdot\,; \hat{\boldsymbol{\lambda}}) - f^*\|_\infty = O_p(r_n)$
  for a rate $r_n \to 0$.
\end{enumerate}
Define the test statistic
\begin{equation}
\label{eq:T_stat}
T_j \;=\;
n \cdot \bigl[\mathrm{CV}(\hat{\boldsymbol{\lambda}}_{-j})
- \mathrm{CV}(\hat{\boldsymbol{\lambda}})\bigr],
\end{equation}
where $\hat{\boldsymbol{\lambda}}_{-j}$ denotes the cross-validated
hyperparameters with $\mu_j$ constrained to zero.
Under $H_0\!: \mu_j^* = 0$, $T_j = o_p(1)$.
Under $H_1\!: \mu_j^* > 0$, $T_j \to \infty$ in probability.
\end{proposition}

\begin{proof}
\emph{The proof is given in Appendix~\ref{app:proofs}.}
\end{proof}

\begin{remark}[Permutation inference]
\label{rem:permutation}
A full distributional theory for $T_j$ under $H_0$ (critical values,
size control, power analysis) requires characterizing the joint
distribution of the cross-validated hyperparameters, which we defer
to future work.
In practice, we implement permutation-based inference:
randomly reassign the structural penalty $C_j$ to pseudo-models
(scrambling the economic content while preserving the penalty's
functional form) and compare $T_j$ to the permutation distribution.
This yields valid $p$-values under exchangeability of the
pseudo-model assignments.
\end{remark}

\begin{remark}[Interpretation]
\label{rem:test_interpretation}
The test~\eqref{eq:test_mu} asks a fundamentally different question
from testing whether a model's pricing errors are zero.
A model may have large pricing errors ($C_j(f^*) > 0$) yet still
receive $\mu_j^* > 0$ if its restrictions are correlated with the
true $f^*$ in a way that improves finite-sample estimation.
Conversely, a correctly specified model ($C_j(f^*) = 0$) may have
$\mu_j^* = 0$ if the RKHS and data are rich enough that the model's
restrictions are redundant.
The test thus measures the \emph{marginal predictive value} of each
theory in the presence of all others.
\end{remark}

\begin{remark}[Relation to model confidence sets]
\label{rem:mcs}
The pairwise comparison underlying~\eqref{eq:T_stat} connects to
the model confidence set (MCS) framework of
\citet*{HansenLundeNason2011}.
One could construct an MCS of theories whose marginal contributions
are statistically indistinguishable:
$\widehat{\mathrm{MCS}}_\alpha = \{j : T_j \leq c_\alpha\}$,
where $c_\alpha$ is calibrated by the permutation distribution.
This provides a principled way to report which theories are
``useful'' for asset-pricing prediction at a given confidence level.
\end{remark}

%======================================================================
\subsection{Summary of the Framework}
\label{sub:summary_framework}
%======================================================================

Table~\ref{tab:framework_summary} collects the main components of the
theory-informed KRR estimator.

\begin{table}[htbp]
\centering
\caption{Summary of the theory-informed KRR framework}
\label{tab:framework_summary}
\small
\begin{tabular}{@{}lll@{}}
\toprule
\textbf{Component} & \textbf{Role} & \textbf{Reference} \\
\midrule
Prediction loss & Fit to observed returns & Eq.~\eqref{eq:tikrr} \\
RKHS penalty & Agnostic statistical regularization & Eq.~\eqref{eq:tikrr} \\
Type~A penalty & Euler equation restrictions (SDF models) & Def.~\ref{def:euler_penalty} \\
Type~B penalty & Production/investment FOC restrictions & Def.~\ref{def:prod_penalty} \\
Type~C penalty & Factor structure restrictions & Def.~\ref{def:factor_penalty} \\
Type~D penalty & Demand-system equilibrium restrictions & Def.~\ref{def:demand_penalty} \\
Representer theorem & Finite-dimensional reduction & Prop.~\ref{prop:representer} \\
First-order conditions & Closed-form for quadratic penalties & Prop.~\ref{prop:foc} \\
Nesting & Recovers KRR, GMM, IPCA, Lasso variants & Prop.~\ref{prop:nesting} \\
CV selection & Grouped penalties, temporal LOYO & Eq.~\eqref{eq:cv} \\
Misspecification & Misspecified theories auto-downweight & Prop.~\ref{prop:misspecification} \\
Inference & Test marginal value of each theory & Prop.~\ref{prop:test_mu} \\
\bottomrule
\end{tabular}
\end{table}
